\section{Sujet n\textsuperscript{o}\,7}
\noindent\begin{minipage}{0.5\linewidth}\footnotesize Office du Baccalauréat du Cameroun\end{minipage}%
\hfill\begin{minipage}{0.45\linewidth}\footnotesize\raggedleft Examen : \textbf{Baccalauréat} · Série : \textbf{C/E}\\
Épreuve : \textbf{Mathématiques} · Durée : \textbf{4\,h} · Coef. : \textbf{7}\end{minipage}
\par\smallskip{\color{onbo}\rule{\linewidth}{1pt}}

\subsection*{Partie A — Évaluation des ressources \hfill (15 points)}

\noindent\textbf{Exercice 1.} \hfill\textit{(3 points)}\\
Dans un centre de santé, une équipe de vaccination est composée de $5$ infirmiers et
$4$ médecins. On forme au hasard une équipe de garde de $3$ personnes.
\begin{enumerate}[label=\arabic*)]
\item Déterminer le nombre total d'équipes de garde possibles. \hfill\textit{(0,75 pt)}
\item Soit $Y$ le nombre de médecins dans l'équipe de garde. Déterminer la loi de
probabilité de $Y$. \hfill\textit{(1,5 pt)}
\item Calculer la probabilité que l'équipe comporte \emph{au moins un} médecin. \hfill\textit{(0,75 pt)}
\end{enumerate}

\smallskip
\noindent\textbf{Exercice 2.} \hfill\textit{(3 points)}\\
On numérote les doses d'un vaccin. Un lot contient un nombre $N$ de doses tel que :
en les regroupant par paquets de $7$, il reste $3$ doses ; en les regroupant par paquets
de $5$, il reste $1$ dose. On suppose de plus que $0\le N<35$.
\begin{enumerate}[label=\arabic*)]
\item Justifier que $5\wedge7=1$ et résoudre dans $\Z$ l'équation $7u-5v=1$. \hfill\textit{(1,5 pt)}
\item Traduire les conditions par un système de congruences et déterminer la valeur de $N$. \hfill\textit{(1,5 pt)}
\end{enumerate}

\smallskip
\noindent\textbf{Exercice 3.} \hfill\textit{(4 points)}\\
La proportion de population immunisée $t$ semaines après le début d'une campagne est
modélisée par $g(t)=1-\mathrm{e}^{-0{,}3t}$, pour $t\ge0$.
\begin{enumerate}[label=\arabic*)]
\item Étudier les variations de $g$ sur $[0,+\infty[$ et déterminer $\displaystyle\lim_{t\to+\infty}g(t)$ ; interpréter. \hfill\textit{(1 pt)}
\item Déterminer, par le calcul, la semaine à partir de laquelle plus de \SI{80}{\percent} de
la population est immunisée. \hfill\textit{(1,5 pt)}
\item À l'aide d'une intégration par parties, calculer
$\displaystyle K=\int_0^{5}t\,\mathrm{e}^{-0{,}3t}\,\dd t$ (valeur exacte puis arrondie à $10^{-2}$). \hfill\textit{(1,5 pt)}
\end{enumerate}

\smallskip
\noindent\textbf{Exercice 4.} \hfill\textit{(5 points)}\\
Le plan complexe est muni d'un repère orthonormé direct $(O\,;\,\vec u,\vec v)$. Trois
centres de vaccination $A$, $B$, $C$ ont pour affixes respectives
$z_A=2+2i$, $z_B=-1+i$ et $z_C=4i$.
\begin{enumerate}[label=\arabic*)]
\item Placer sommairement $A$, $B$, $C$ et calculer les distances $AB$, $AC$ et $BC$. \hfill\textit{(1,5 pt)}
\item En déduire la nature précise du triangle $ABC$. \hfill\textit{(1 pt)}
\item Soit $S$ la similitude directe de centre $A$, de rapport $\frac{\sqrt2}{2}$ et
d'angle $-\frac\pi2$. Donner l'écriture complexe de $S$. \hfill\textit{(1,5 pt)}
\item Déterminer l'affixe de l'image $B'$ du point $B$ par $S$. \hfill\textit{(1 pt)}
\end{enumerate}

\subsection*{Partie B — Évaluation des compétences \hfill (5 points)}

\noindent\textit{Le district de santé de Maroua organise une campagne de vaccination. Tu es
sollicité, en tant qu'élève de Terminale~C, pour appuyer trois décisions. Le premier jour,
$500$ personnes sont vaccinées ; grâce à la sensibilisation, ce nombre augmente de
\SI{15}{\percent} chaque jour. Par ailleurs, un test de dépistage utilisé avant la vaccination
est fiable : il est positif pour \SI{96}{\percent} des personnes réellement malades et
négatif pour \SI{98}{\percent} des personnes saines ; on sait que \SI{2}{\percent} de la
population est réellement malade. Enfin, on prélève au hasard $6$ flacons d'un lot où
\SI{3}{\percent} des flacons sont périmés.}

\begin{enumerate}[label=\textbf{Tâche \arabic*.},leftmargin=2.4em]
\item Déterminer le jour à partir duquel le nombre de personnes vaccinées dans la journée
dépassera $1500$. \hfill\textit{(1,5 pt)}
\item Une personne est testée positive. Déterminer la probabilité qu'elle soit
réellement malade (arrondie à $10^{-2}$). \hfill\textit{(1,5 pt)}
\item Déterminer la probabilité qu'\emph{aucun} des $6$ flacons prélevés ne soit périmé. \hfill\textit{(1,5 pt)}
\end{enumerate}
\par\smallskip{\footnotesize\itshape Présentation générale : 0,5 point.}

% ════════════════════════════════════════════════════
\corrige{du sujet 7}

\noindent\textbf{Partie A — Exercice 1.}
1) $\binom{9}{3}=84$ équipes possibles.
2) $Y\in\{0,1,2,3\}$, $P(Y{=}k)=\dfrac{\binom4k\binom5{3-k}}{84}$ :
$P(0)=\frac{10}{84}=\frac5{42}$, $P(1)=\frac{40}{84}=\frac{10}{21}$,
$P(2)=\frac{30}{84}=\frac5{14}$, $P(3)=\frac{4}{84}=\frac1{21}$.
(Vérification : $10+40+30+4=84$.)
3) $P(Y\ge1)=1-P(Y{=}0)=1-\frac5{42}=\frac{37}{42}$.

\noindent\textbf{Exercice 2.}
1) $7=1\times5+2$, $5=2\times2+1$ : le dernier reste non nul est $1$, donc $5\wedge7=1$.
De $5=2\cdot2+1$ et $2=7-5$ : $1=5-2(7-5)=3\times5-2\times7$, soit $7\times(-2)-5\times(-3)=1$ ;
une solution particulière est $(u,v)=(-2,-3)$, et la solution générale
$u=-2+5k,\ v=-3+7k$, $k\in\Z$.
2) Système : $N\equiv3\ [7]$ et $N\equiv1\ [5]$. Les multiples de $7$ valant $3$ modulo $7$\dots
on cherche $N$ : $N\equiv3\,[7]$ donne $N\in\{3,10,17,24,31\}$ ; parmi eux, $N\equiv1\,[5]$
impose $N\equiv1\,[5]$ : $3,10,17,24,31\to$ restes mod $5$ : $3,0,2,4,1$. Seul $N=31$ convient.
Donc $\boxed{N=31}$ (et $31=4\times7+3=6\times5+1$).

\noindent\textbf{Exercice 3.}
1) $g'(t)=0{,}3\,\mathrm{e}^{-0{,}3t}>0$ : $g$ est strictement croissante. $g(0)=0$ et
$\lim_{t\to+\infty}g(t)=1-0=1$ : à terme, la quasi-totalité de la population est immunisée.
2) $g(t)>0{,}8\iff\mathrm{e}^{-0{,}3t}<0{,}2\iff-0{,}3t<\ln0{,}2\iff t>\dfrac{-\ln0{,}2}{0{,}3}=\dfrac{\ln5}{0{,}3}\approx5{,}36$.
À partir de la \textbf{$6^{\mathrm e}$ semaine}.
3) IPP ($w=t$, $w'=1$, $h'=\mathrm{e}^{-0{,}3t}$, $h=-\frac1{0{,}3}\mathrm{e}^{-0{,}3t}=-\frac{10}{3}\mathrm{e}^{-0{,}3t}$) :
\[K=\Big[-\tfrac{10}{3}t\,\mathrm{e}^{-0{,}3t}\Big]_0^5+\tfrac{10}{3}\int_0^5\mathrm{e}^{-0{,}3t}\dd t
=-\tfrac{50}{3}\mathrm{e}^{-1{,}5}+\tfrac{10}{3}\Big[-\tfrac{10}{3}\mathrm{e}^{-0{,}3t}\Big]_0^5.\]
Soit $K=-\frac{50}{3}\mathrm{e}^{-1{,}5}-\frac{100}{9}(\mathrm{e}^{-1{,}5}-1)
=\frac{100}{9}-\Big(\frac{50}{3}+\frac{100}{9}\Big)\mathrm{e}^{-1{,}5}
=\frac{100}{9}-\frac{250}{9}\mathrm{e}^{-1{,}5}\approx11{,}11-6{,}20\approx\mathbf{4{,}91}$.

\noindent\textbf{Exercice 4.}
1) $AB=|z_B-z_A|=|-3-i|=\sqrt{9+1}=\sqrt{10}$ ;
$AC=|z_C-z_A|=|-2+2i|=\sqrt{4+4}=2\sqrt2$ ;
$BC=|z_C-z_B|=|1+3i|=\sqrt{1+9}=\sqrt{10}$.
2) $AB=BC=\sqrt{10}$ : le triangle $ABC$ est \textbf{isocèle} en $B$.
(De plus $AB^2+\,?$ : $AC^2=8$, $AB^2=BC^2=10$ ; il n'est pas rectangle.)
3) $S(z)=a(z-z_A)+z_A$ avec $a=\frac{\sqrt2}{2}\mathrm{e}^{-i\pi/2}=\frac{\sqrt2}{2}(-i)=-\frac{\sqrt2}{2}i$.
Donc $S(z)=-\frac{\sqrt2}{2}i\,(z-(2+2i))+2+2i$.
4) $z_B-z_A=-3-i$. $a(z_B-z_A)=-\frac{\sqrt2}{2}i(-3-i)=-\frac{\sqrt2}{2}(-3i-i^2)
=-\frac{\sqrt2}{2}(1-3i)=-\frac{\sqrt2}{2}+\frac{3\sqrt2}{2}i$.
Puis $z_{B'}=a(z_B-z_A)+z_A=\big(2-\tfrac{\sqrt2}{2}\big)+\big(2+\tfrac{3\sqrt2}{2}\big)i$.

\smallskip
\noindent\textbf{Partie B — Situation.}
\emph{Tâche 1.} Nombre de personnes le jour $n$ : $v_n=500\times1{,}15^{\,n-1}$. On résout
$v_n>1500\iff1{,}15^{\,n-1}>3\iff n-1>\dfrac{\ln3}{\ln1{,}15}\approx\dfrac{1{,}0986}{0{,}1398}\approx7{,}86$,
soit $n>8{,}86$ : à partir du \textbf{$9^{\mathrm e}$ jour} ($v_9\approx500\times1{,}15^8\approx1529$).

\emph{Tâche 2.} Notons $M$ « malade », $T^+$ « test positif ». $P(M)=0{,}02$,
$P(T^+\mid M)=0{,}96$, $P(T^+\mid\bar M)=1-0{,}98=0{,}02$.
$P(T^+)=0{,}02\times0{,}96+0{,}98\times0{,}02=0{,}0192+0{,}0196=0{,}0388$.
$P(M\mid T^+)=\dfrac{0{,}0192}{0{,}0388}\approx\mathbf{0{,}49}$ (environ \SI{49}{\percent}).

\emph{Tâche 3.} $P(\text{aucun périmé})=(0{,}97)^6\approx\mathbf{0{,}833}$, soit environ \SI{83,3}{\percent}.
